\section{The theorem and its constructive meaning}
\label{sec:statement}

Brun's theorem asserts that the reciprocals in all twin-prime pairs have
a finite sum. We index a pair by its smaller member. Write
\[
 \Twin(p)\quad\Longleftrightarrow\quad
 p\text{ and }p+2\text{ are prime},
\]
and define, for natural numbers \(X,N\),
\[
 T(X)=\#\{p:1\le p\le X,\ \Twin(p)\},\qquad
 B(N)=\sum_{\substack{1\le p<N\\\Twin(p)}}
       \left(\frac1p+\frac1{p+2}\right).
\]
The convention deliberately counts \(5\) twice: once in \((3,5)\) and once
in \((5,7)\). Thus \(B(4)=8/15\) and \(B(6)=92/105\).
The endpoint in \(T\) is inclusive; the endpoint in \(B\) is strict.
The condition \(p+2\le X\) is \emph{not} imposed in the definition of \(T(X)\).

\begin{theorem}[Finite quantitative Brun theorem]
\label{thm:main}
The following assertions have proofs in intuitionistic primitive recursive
arithmetic, when rational notation is read through positive-denominator
integer codes.
\begin{enumerate}[label=(\roman*)]
\item For every \(k\ge0\),
\begin{equation}
 k^2T(2^k)\le 1400\,2^k\le 2^{11}2^k.
 \label{eq:main-count}
\end{equation}
\item For all \(m\ge1\) and \(L\ge0\),
\begin{equation}
 \sum_{\substack{2^m\le p<L\\\Twin(p)}}
 \left(\frac1p+\frac1{p+2}\right)
 \le \frac{2^{13}}m.
 \label{eq:main-tail}
\end{equation}
\item Define \(N(q)=2^{\,2^{q+13}}\). For all \(q,M,L\in\N\),
\begin{equation}
 M,L\ge N(q)\quad\Longrightarrow\quad
 \lvert B(M)-B(L)\rvert\le 2^{-q}.
 \label{eq:main-cauchy}
\end{equation}
\end{enumerate}
\end{theorem}

The constants are intentionally generous. They simplify the arithmetic
proof and give a completely specified modulus; they are not intended for
efficient numerical evaluation of Brun's constant. In particular, even the
cutoff at \(q=0\) is \(2^{8192}\).

\subsection{A statement literally in the language of iPRA}

There are no real-number variables in first-order iPRA. Accordingly, the
primary conclusion is a theorem about finite integer computations. Here is
one exact version.

Let \(\chi(p)\in\{0,1\}\) be the primitive recursive characteristic function
of \(\Twin(p)\). Primality requires \(p\ge2\) and is tested by trial division
over the explicitly bounded range \(2,\ldots,p-1\). Put
\[
 D(V)=(V+1)!,
\]
and define the integer
\begin{equation}
 R(U,V)=
 \sum_{U\le p<V}\chi(p)
 \left(D(V)\mathbin{\mathrm{div}}p+
       D(V)\mathbin{\mathrm{div}}(p+2)\right).
 \label{eq:integer-tail}
\end{equation}
Use total division, with a fixed value such as \(0\) at divisor \(0\).
The terms at \(p=0,1\) have coefficient \(0\). Every active denominator is
positive and at most \(V+1\), and hence divides \(D(V)\).
It follows exactly, without rounding, that for \(U\le V\),
\begin{equation}
 B(V)-B(U)=\frac{R(U,V)}{D(V)}.
 \label{eq:tail-code}
\end{equation}

The arithmetical theorem to be proved is the universal closure of
\begin{equation}
 \boxed{\quad
 N(q)\le U\ \wedge\ U\le V
 \quad\Longrightarrow\quad
 2^qR(U,V)\le D(V).
 \quad}
 \label{eq:ipra-target}
\end{equation}
All symbols in its matrix denote primitive recursive functions or
relations with primitive recursive characteristic functions. In particular,
its matrix is a quantifier-free formula after replacing order and bounded
tests by those characteristic functions. Formula~\eqref{eq:ipra-target}
therefore does not mean an informal quantification over infinite sums.

The order on natural numbers is decidable. Applying
\eqref{eq:ipra-target} to the smaller and larger of \(M,L\) gives
\eqref{eq:main-cauchy}. Conversely, the finite rational tail estimate proves
\eqref{eq:ipra-target} by multiplying by the positive integer \(2^qD(V)\).

\subsection{What convergence means here}

The sequence
\[
 b(q)=B\!\left(2^{\,2^{q+13}}\right)
\]
is a specified primitive recursive rational name. For \(r\ge q\),
\(\lvert b(r)-b(q)\rvert\le2^{-q}\).
In a constructive development of Cauchy reals, this name represents Brun's
constant. Its equivalence to the original sequence of partial sums follows
from the estimate for \emph{all} \(M,L\) beyond the cutoff, not merely for
the selected subsequence.

If the convention for a Cauchy modulus requires a strict inequality
\(<2^{-q}\), replace \(N(q)\) by \(N(q+1)\).
The proof never needs to decide whether there are finitely many or infinitely
many twin-prime pairs. A convergent series of nonnegative terms is compatible
with either possibility.

\subsection{Proof architecture}

The classical Selberg sieve provides the organizing idea
\cite{kedlaya,tau}, but we will prove the particular finite identity needed.
The steps are:
\begin{enumerate}
\item Count the forbidden residues of \(n(n+2)\) modulo squarefree integers.
\item Construct explicit rational coefficients \(\lambda_d\) whose squared
divisor sum majorizes the indicator of surviving the sieve.
\item Evaluate its main quadratic form as \(1/G(z)\), and bound the total
counting error by \(z^4\).
\item Prove \(G(s^2)\ge H_s^2/3\) by a finite prime-factor grouping.
\item Choose dyadic parameters and telescope a finite rational majorant
for the reciprocal tail.
\end{enumerate}
Sections~\ref{sec:finite}--\ref{sec:convergence} carry out these steps.
Section~\ref{sec:ipra} checks their induction and coding requirements.
